% Equations from chunk_024 (pages 231-240 / internal pages 208-217)
% Chapter 10: The Net Social Wage in Turkey, 1980-2019
% Authors: Y. Karabacak and E. A. Tonak

% Equation 1: Net Social Wage Definition
% Page 208 (internal numbering)
% Context: Basic definition of net social wage (NSW)
\begin{equation}
\pm NSW = B - T
\end{equation}
% Where:
% NSW = Net Social Wage
% B = Benefits (net gains and benefits in kind or cash)
% T = Taxes collected by the state from wages and salaries

% Equation 2: Public Expenditure Benefiting Working Class
% Page 213 (internal numbering)
% Context: Calculation of benefits to working class from public expenditures
\begin{equation}
B = (ls \times B_2) + B_3
\end{equation}
% Where:
% B = Total benefits to working class
% ls = Labour share (wages and salaries/personal income)
% B_2 = Public expenditures that benefit all (fuel/energy, transport, environmental protection, housing, healthcare, recreation, culture, religious services, education)
% B_3 = Public expenditures benefiting working class entirely (social security, social assistance, general labour affairs, unemployment benefits)

% Equation 3: Taxes Paid by Working Class
% Page 215 (internal numbering)
% Context: Calculation of total taxes paid by working class
\begin{equation}
T = T_1 + (ls \times T_2)
\end{equation}
% Where:
% T = Total taxes paid by working class
% T_1 = Taxes paid entirely by working population (social security premia, unemployment insurance fund deductions)
% ls = Labour share coefficient
% T_2 = Taxes paid by all sections of population (personal income tax, VAT, consumption taxes, etc.)
% Note: T_3 (taxes paid by those outside working class) not included in this equation

% Additional notation used throughout:
% B_1 = Public expenditures with zero labour share (general public services, defence, public order and security, economic affairs)
% T_3 = Taxes paid by non-working class (corporate income tax, inheritance tax, real estate tax, etc.)
% GDP = Gross Domestic Product
% GNP = Gross National Product
